\section{Introduction}\label{sec:intro}

Many optimisation problems are \emph{integer linear programs} (ILPs): a linear objective and linear constraints over integer variables. Solving an ILP is NP-hard in general \cite{karp1972}, whereas a \emph{linear program} (LP) can be solved in polynomial time \cite{khachiyan1979}. A standard approximation recipe exploits this gap: write the problem as an ILP, drop integrality to get the \emph{LP relaxation}, solve it and round. For a minimisation problem the relaxation can only lower the optimum, $\optlp \le \optilp$, so rounding usually costs something. For vertex cover, rounding at $\tfrac12$ only guarantees a factor of~2 \cite{hochbaum1982}.

Sometimes rounding costs nothing. The \emph{vertex cover LP} (VC LP) of a graph $G=(V,E)$ is
\[
  \min \sum_{v\in V} x_v
  \quad\text{s.t.}\quad x_u + x_v \ge 1 \ \text{ for every } \{u,v\}\in E,
  \qquad 0 \le x \le 1 .
\]
On the triangle $C_3$ the all-$\tfrac12$ point is optimal with value $1.5$ (summing the edge constraints gives $2\sum_v x_v\ge3$), but every cover needs two vertices. On the 4-cycle $C_4$ the LP optimum is~$2$, and simplex returns an integral cover (Figure~\ref{fig:running}). This report asks three questions:

\begin{quote}\itshape
  When does solving an LP give an integer optimum? How do we recognise such problems? Does the LP algorithm we use matter?
\end{quote}

\begin{figure}[H]
  \centering
  % Figure~\ref{fig:running}: VC LP on C_3 and C_4.
% Input by sections/01-introduction.tex.
\begin{tikzpicture}[
    vtx/.style={circle, draw, minimum size=6.5mm, inner sep=0pt, font=\small},
    lbl/.style={font=\small}]
  \begin{scope}
    \node[vtx, fill=warn!18] (a) at (90:0.85)  {$\tfrac12$};
    \node[vtx, fill=warn!18] (b) at (210:0.85) {$\tfrac12$};
    \node[vtx, fill=warn!18] (c) at (330:0.85) {$\tfrac12$};
    \draw (a) -- (b) -- (c) -- (a);
    \node[lbl, anchor=west] at (1.25,0) {$C_3$: \textcolor{warn}{$\optlp = 1.5$}, $\optilp = 2$};
  \end{scope}
  \begin{scope}[xshift=8.2cm]
    \node[vtx, fill=good!22] (p) at (-0.6, 0.6) {$1$};
    \node[vtx]               (q) at ( 0.6, 0.6) {$0$};
    \node[vtx, fill=good!22] (r) at ( 0.6,-0.6) {$1$};
    \node[vtx]               (s) at (-0.6,-0.6) {$0$};
    \draw (p) -- (q) -- (r) -- (s) -- (p);
    \node[lbl, anchor=west] at (1.05,0) {$C_4$: \textcolor{good}{$\optlp = \optilp = 2$}};
  \end{scope}
\end{tikzpicture}%

  \caption{The running example, labelled with optimal LP solutions $x_v$. On $C_3$ the LP optimum is fractional and below the integer optimum; on $C_4$ simplex returns an integral optimal cover.}
  \label{fig:running}
\end{figure}

The answer centres on \emph{total unimodularity} (TU): every square submatrix of the constraint matrix $A$ has determinant $-1$, $0$ or~$1$. If $A$ is TU and the right-hand side $b$ is integral, every vertex of the feasible region is integral \cite{hoffman1956}, so the LP has an integral optimum. The matrix of $C_4$ is TU because $C_4$ is bipartite; that of $C_3$ is not, since its determinant is~$2$.

\paragraph{Scope and outline.} Figure~\ref{fig:roadmap} maps the questions to our sections and experiments. We prove the direction of the Hoffman--Kruskal theorem that we use, derive the main TU families, compare three ways to recognise TU, and study how LP algorithms behave on these LPs. We state Seymour's theorem \cite{seymour1980} without proof and use an existing implementation of its test, and only mention total dual integrality and the non-bipartite matching polytope. All code, data and raw results are in our repository.\footnote{\url{https://github.com/feilinlp/linear-programming-total-unimodularity}} Every number in this report was produced there, and each experiment reruns with one command (Appendix~\ref{app:repro}).

\begin{figure}[H]
  \centering
  % Figure~\ref{fig:roadmap}: questions, theory and experiments.
% Input by sections/01-introduction.tex.
\begin{tikzpicture}[font=\small,
    q/.style={rectangle, rounded corners=2pt, fill=accent!12, draw=accent!60,
              text width=4.85cm, align=center, inner sep=3pt, minimum height=7mm},
    th/.style={rectangle, draw=black!25, text width=4.85cm, align=flush left,
               inner sep=3pt, minimum height=13mm},
    ex/.style={rectangle, rounded corners=2pt, fill=black!6, draw=black!25,
               text width=4.85cm, align=flush left, inner sep=3pt, minimum height=10.5mm}]
  \foreach \x/\name/\qq/\theory/\exper in {
      0/1/{\textbf{When} is the LP exact?}/{\S\ref{sec:tu}: Hoffman--Kruskal \cite{hoffman1956}, TU families, half-integrality \cite{nemhauser1975}}/{\textbf{E1}: is TU necessary?},
      5.35/2/{\textbf{How} do we recognise TU?}/{\S\ref{sec:recognise}: brute force, Ghouila-Houri \cite{ghouilahouri1962}, Seymour \cite{seymour1980}}/{\textbf{E2}: how costly is recognition?},
      10.7/3/{Does the \textbf{algorithm} matter?}/{\S\ref{sec:solve}: simplex \cite{dantzig1951}, ellipsoid \cite{khachiyan1979}, IPM \cite{karmarkar1984}, crossover \cite{megiddo1991}}/{\textbf{E3}: which optimal point?\\ \textbf{E4}: LP or specialised code?}} {
    \node[q]  (q\name) at (\x, 0) {\qq};
    \node[th, below=1mm of q\name.south, anchor=north] (t\name) {\theory};
    \node[ex, below=1mm of t\name.south, anchor=north] (e\name) {\exper};
  }
\end{tikzpicture}%

  \caption{Roadmap: the three questions, the theory and main papers that answer each, and the experiments that test them.}
  \label{fig:roadmap}
\end{figure}
